\documentclass{article}
\usepackage{graphicx} % Required for inserting images
\usepackage{parskip} % Removes paragraph indents

\title{\textbf{Determining the Coefficient of Static and Kinetic Friction for Difference Surfaces}}
\author{Ned Hillyard}
\date{September 2026}

\begin{document}

\maketitle

\section*{Abstract}
save till last 

\section*{Introduction}

Friction is the resistive force between two surfaces that opposes the motion of the surfaces. It's an important physical phenomena because it's the resisting force that allows us to walk around, grab and hold things and stop cars via braking. When an object is in direct contact with a horizontal surface and a force is applied, the maximum friction force that can be applied before motion is induced is the static friction. The coefficient of static friction is the ratio between the maximum static friction force and the normal force.

\begin{equation}
    \mu_s = \frac{F_s}{F_N} = \frac{F_{max}}{mg}
\end{equation}

Once the object is in motion and at a constant velocity, the opposing force becomes kinetic friction. The coefficient of kinetic friction is similarly defined as the ratio between the kinetic frictional force and the normal force.

\begin{equation}
    \mu_k = \frac{F_k}{F_N} = \frac{F}{mg}
\end{equation}

And these two coefficients $\mu_s$ and $\mu_k$ characterise how strongly the two surfaces interact. The aim of this experiment was to measure these coefficients of friction for several surfaces (cork, felt and plastic), and to examine how exactly friction depends on normal force, surface area and velocity.

\section*{Methods}

\section*{results}

\section*{Discussion}

\section*{Conclusion}

\section*{References}

\end{document}
